% English version of paper/ja_v2/09_結論.md (2026-11-17). The closing suggestion sentence was removed at the researcher's request (2026-11-17): the conclusion ends with the proposal and its result.
\section{Conclusion}\label{sec:conclusion}

We treated multimodal emotion recognition on a new corpus with a few labels per emotion as semi-supervised domain adaptation and answered three questions with linear classifiers on frozen pretrained encoders. With one source corpus, audio and visual answer bias was almost removed by subtracting feature means without labels; with all other groups as sources, the remaining error was rarely reduced by same-modality unlabelled utterances and was reduced by pseudo-labels from classifiers including other modalities. When AFEW, IEMOCAP (improvised), MOSEI and MELD each served as the new corpus, distillation on these pseudo-labels gave the audio classifier 1.4--4.1\pt\ higher UAR than before retraining in all 12 conditions with 5, 10 and 25 labels per emotion (point estimates). Across 360 post hoc pseudo-label variants per K (K = 5, 25; pseudo-labels made only of the retrained classifier's own predictions and duplicates of other providers excluded), the rank correlation between the advantage of the pseudo-labels over the classifier before retraining and the UAR difference was 0.80 and 0.83, and no variant with a negative advantage rose. Because the strongest modality differed by corpus and emotion, we proposed late fusion that measures strengths on the few target labels and per-emotion profiles on the sources, with hyperparameters selected on source-corpus episodes. Evaluated on the four corpora used to design it, the method had higher UAR than the product of the three classifiers' probabilities in 11 of 12 conditions in point estimates, by 2.1--3.5\pt\ on AFEW.
